\odiseachapter{homero}{Homero}\label{cap:homero}

Si es que hubo un Homero, claro. Infinitas son las disputas de los magos y más infinitos aún los debates sobre si Homero fue una sola persona o un colectivo que edita y unifica una rica tradición oral anterior\footnote{Una introducción a la cuestión homérica puede encontrarse en Gregory Nagy, \emph{Homeric Questions}: \url{https://chs.harvard.edu/book/nagy-gregory-homeric-questions/}}. También se discute vivamente si la \emph{Ilíada} y la \emph{Odisea} pertenecen al mismo autor.

La noción de que ambos poemas épicos se derivan de una amplia ---y muy antigua--- tradición oral tiene ya cien años. El personaje más importante asociado a la idea es Milman Parry (1902--1935), quien ---nobleza obliga--- tuvo un final prematuro y absurdo, al dispararse por accidente su propio revólver. Imposible no pensar en los caprichosos dioses.

La idea de Parry resulta muy convincente para cualquiera que haya visto en acción a los bertsolaris de Euskadi o se haya deleitado con las «peleas de gallos», en las que jóvenes raperos se dan bofetadas verbales, inventando sus versos al vuelo\footnote{Mi hijo es muy aficionado a estas peleas de gallos y hace poco escuchábamos una en la que se esgrimían con descaro estas líneas: «Tú usas la violencia / yo soy inteligente / ¡y te hago bullying... mentalmente!»}. Unos y otros improvisan, sí, pero basándose en patrones bien establecidos y memorizados que les permiten ejecutar en tiempo real variaciones sobre un tema y les facilitan el ritmo y la rima. Es algo parecido a lo que se hace en jazz, donde la inagotable creatividad de los músicos se ancla en un tronco musical muy robusto.

Parry argumentó exactamente eso, mostrando que tanto la \emph{Ilíada} como la \emph{Odisea} usan fórmulas orales, temas fijos y escenas más o menos sincopadas cuyo objetivo es permitir la improvisación al vuelo. A lo largo de estas páginas hemos visto muchas de estas fórmulas. Una de mis favoritas se usa para nombrar a Ulises:

\begin{versocita}
Sangre Laercia, Odiseo milmañas, hijo del cielo
\end{versocita}

Ya vimos que las fórmulas de cortesía son idénticas. Néstor usa la misma para agasajar a Telémaco que Polifemo para burlarse de los aqueos. Los ejemplos abundan. Atenea es ojiglauca, Zeus es cabrillante, Helena es mantodevuelo y cabellilinda, Circe beltrenza y vocihumana, Penélope milpensamientos. Todo un armazón de trucos nemotécnicos y expresiones prefabricadas disponible para que el bardo pueda memorizar e improvisar. Dicho sea de paso, los versos de la traducción de Tobal introducen además la rima, lo que facilita aún más el trabajo de la memoria y hace más placentero declamar en voz alta.

Parry llamó a estos patrones fijos, precisamente, «fórmulas». Ya hemos visto algunas; ahí van otras: \gr{ῥοδοδάκτυλος Ἠώς} (rhododáktylos Ēṓs) o «Aurora de rosados dedos»; \gr{ἔπεα πτερόεντα} (épea pteróenta) o «aladas palabras»; \gr{πολύμητις Ὀδυσσεύς} (polýmētis Odysseús), que Tobal traduce como Odiseo mientisagaz o milastucias; y \gr{ἕρκος ὀδόντων} (hérkos odóntōn) o «el cerco de los dientes».

La fértil intuición de Parry ha sido extendida y cuantificada a lo largo de las décadas, incluyendo estudios estadísticos que muestran que la densidad de «fórmulas» es muy superior a la que se encuentra en autores clásicos posteriores, como Eurípides o Virgilio.

¿Qué demuestra entonces el análisis de Parry? En el debate anterior a su trabajo, un argumento repetido a favor del concepto de los múltiples autores era el aparente batiburrillo de estilos y capas dialectales, que sugería que el poema había sido compuesto por varios bardos desparramados en distintas épocas. Milman mientisagaz demuestra que repeticiones, dobladillos, cantinelas y remiendos eran en realidad parte de la técnica de composición. Homero podría seguir siendo un solo autor, posiblemente ---pero no necesariamente--- literato, que recoge la rica tradición heredada en forma de relatos fragmentados y le da coherencia, creando una épica unitaria.

Homero sería entonces la voz en cuyo eco resuenan muchas voces, un autor único y a la vez plural. Puede no ser la única interpretación válida, pero debo confesar aquí que es la que más me gusta.

Aunque todavía me gusta más la traviesa alternativa que planteó el célebre novelista Samuel Butler ---autor de \emph{Erewhon} y \emph{The Way of All Flesh}--- en 1897, hace la friolera de (casi) 130 años. En \emph{The Authoress of the Odyssey}, Butler sostenía que la \emph{Odisea} había sido escrita por una muchacha de Trápani, en Sicilia, que se autorretrató como Nausícaa. No me queda claro si Butler se tomó del todo en serio su hipótesis ---aunque la defendió a capa y espada durante la última década de su vida--- o si fue una de estas bromas cultas que, como el vino dulce, acaban por subirse a la cabeza. En todo caso, la idea inspiró a Robert Graves la divertida novela \emph{La hija de Homero}.

La novela ---en este punto los temerosos del \emph{spoiler} harían bien en saltarse el capítulo--- relata las aventuras de Nausícaa ---cuyo nombre, se nos informa, es «la que quema las naves»---, princesa de la casa real de los élimos, en
Drepanum (la moderna Trápani), en Sicilia. El argumento involucra una conspiración de las otras casas nobles contra su familia cuyo objetivo es eliminar a su padre y hermanos y hacerse con el poder. La historia arranca con la desaparición del hermano mayor de Nausícaa, el príncipe Laodamante, durante un viaje cuya motivación es perfectamente insensata, a saber, buscar un collar de ámbar para su desabrida esposa Ctímene. Al cabo de un año de ausencia, su padre, el rey Alfides, decide ir a buscarle y se embarca rumbo a la arenosa Pilos, donde espera recabar noticias suyas. Quedan en la casa las mujeres ---la princesa Nausícaa y su madre---, un hijo todavía adolescente (Clitoneo) y Méntor, el anciano tío materno que ocupa el cargo de regente.

El rey Alfides es un personaje curioso. Amigo de echar sermones, honesto pero agarrado, justo pero solemne, noble pero propenso a irritar a sus pares con su pedantería. Los principales de las casas nobles lo aguantan a duras penas y aprovechan su ausencia para intentar un golpe de Estado en varias fases, la primera de las cuales involucra invadirle la casa con sus jóvenes vástagos, con la excusa de solicitar la mano de la princesa.

Es decir, Robert Graves recrea las circunstancias que llevan a los confiados y altaneros pretendientes a invadir el palacio del rey ausente. La deliciosa técnica de don Roberto se basa en subvertir continuamente los hechos. Sigue habiendo un rey ausente, que tiene algunas de las características de Odiseo ---no le falta astucia ni inteligencia, por ejemplo--- pero se nos presenta como mucho más limitado y mucho más humano. Hay también un hermano mayor ausente y a la vez uno menor que ayuda a la princesa ---Telémaco desdoblándose---, la banda de gorrones es la misma de la \emph{Odisea}, pero ahora no persiguen a Penélope sino a una princesa casadera. La épica se ha convertido en un barullo doméstico y el prisma que el escritor aplica tiene la virtud de hacer la historia más cercana y reconocible.

Nausícaa, por su parte, no solo está decidida a no casarse con ninguno de los pretendientes, sino que trama su ruina ---justificadamente, ya que, en realidad, actúa en defensa propia, los jóvenes están menos interesados en casarse con ella que en asesinar a su padre y su hermano y quedarse con su hacienda---. Sin apenas aliados, ya que solo cuenta con el adolescente Clitoneo y el anciano Méntor, un día, la fortuna ---o la diosa Atenea de la que Nausícaa se sirve con mucha astucia y cierta falta de escrúpulos--- depara que encuentre a un náufrago solitario, mientras ella y sus criadas hacen la colada. El náufrago en cuestión dice llamarse Etón y ser cretense. Gracias a él y con la ayuda de su hermano menor y los criados Eumeo y Filetio, Nausícaa prepara la prueba del arco y ejecuta la matanza de los pretendientes.

Entre tanto, la novela hace aparecer al bardo Demódoco y a su pupilo Femio. Demódoco recrea en sus cantos el \emph{nóstos} de diferentes príncipes aqueos, incluyendo el de Odiseo y Graves ofrece una versión blasfema y desternillante del regreso a Ítaca, en el que Ulises se encuentra con que Penélope, lejos de serle fiel, se acuesta con nada menos que cincuenta pretendientes, una proeza que solo puede ejecutar gracias a la divinal energía que le proporciona el cinturón mágico de Afrodita. Pero a don Roberto no le basta con transformar a la castísima Penélope en la agente de Afrodita en la tierra. Odiseo se pasa siete años varado en una isla, pero sin Calipso. Circe lo trata a puntapiés y en general va de mal en peor hasta llegar a Ítaca. Recuerdo muy bien mi primera lectura de la divertidísima novela ---tendría veintitantos años--- y la felicidad con que cerré filas con el vengativo escritor inglés, que ---no me cupo duda alguna--- le estaba haciendo pagar a Ulises sus desmanes con la misma saña con que yo planeaba hacerlo algún día.

Por si no fueran pocas vueltas de tuerca, la tesis final de la novela es que Nausícaa se sirve del bardo Femio para difundir su propia versión del poema, que es la que llega a nuestros días. Y todo esto me sirve para resumir el punto de vista sobre el que llevo insistiendo a lo largo de estas páginas:

\emph{El clásico existe para que las generaciones posteriores lo reescriban, lo reimaginen, lo reinterpreten.}

Graves hace las tres cosas.

Para empezar, le da la vuelta a la narrativa, transformando el poema en pura ficción. No olvidemos que la \emph{Odisea} se presenta, formalmente, casi como un relato histórico, la descripción literal de las aventuras de uno de los héroes de Troya en su regreso al hogar. En cambio, en \emph{La hija de Homero}, la «realidad», mucho más prosaica, es una conspiración palaciega, que la imaginación de la princesa convierte en épica.

Segundo, reinterpreta los personajes de manera a la vez sutil y atrevida. Nausícaa pasa de ser la princesa a la que Odiseo relata sus aventuras a ser la poeta que las inventa. El cretense Etón es a la vez un trasunto y una subversión de Ulises. Veamos por qué. En el canto XIX, Penélope baja de sus aposentos, se encuentra a Ulises, disfrazado de mendigo, y, tras un breve enfrentamiento con la desdichada criada Melanto, le pregunta:


\begin{versocita}
«Es lo primero, mi huésped, por ti preguntar y que sepa:\\
¿quién y de qué gentes tú? ¿Do la casa y ciudad que te esperan?».
\end{versocita}

A lo que su disfrazado marido responde ---tras unas cuantas florituras---, diciéndole ser cretense:

\begin{versocita}
«Hay una tierra, Creta, en medio del mar vinulento,\\
bella y feraz, de aguas toda a redor, que es de hombres asiento,\\
muchísimos, innumerables, noventa ciudades le cuento.\\
Las lenguas de unos y otros allí se confunden: aqueos\\
hay, y cidones, y eteócretes cordialtaneros,\\
y celestiales pelasgos, y dorios del trino de pueblos».
\end{versocita}

y, ¡sorpresa!, asegurando que su nombre es Etón.

¿Qué hace la hija de Homero? Le da la vuelta al argumento. El personaje «real» ---un joven, bello y zalamero cretense, que sin embargo resulta ser todo un asesino a la hora de la verdad--- es el cretense Etón, y Odiseo una leyenda que se inspira, en parte, en él. Lo más notable del truco de magia que practica Graves es que la escena en la que el náufrago y la princesa se encuentran recrea, casi hasta la coma, los versos épicos, y a la vez se las compone para transmitirnos algo radicalmente diferente. Aparece cada detalle, minuciosamente observado. La alegre tropa de doncellas lavando la ropa y jugando a la vez, el encuentro entre la princesa y el náufrago ---sin escamotear el hecho de que Etón, como Ulises, a pesar de ser muy robusto, es algo patizambo---, el desenlace posterior, que sigue el guion de la \emph{Odisea} al pie de la letra... Hasta que Nausícaa cambia de opinión y lo pone patas arriba, con una naturalidad que deja al lector boquiabierto.

Pero la subversión del argumento, hasta cierto punto, es lo de menos. Lo verdaderamente subversivo es la forma en que el peso de la historia recae en la joven, que no necesita que venga ningún marido a rescatarla de los pretendientes y a salvar su casa, sino que se las ingenia para emplear al héroe de turno en su servicio.

Finalmente, Graves cambia el registro. Homero es dramático siempre, cínico a menudo, trágico muchas veces y casi nunca nos hace reír. En \emph{La hija de Homero}, la feliz lectora, el complacido lector, no pueden dejar de hacerlo. Nausícaa es alegre, vivaz, astuta, descarada, rápida de reflejos, generosa, valiente... Tiene todas las virtudes de Odiseo y muy pocos de sus vicios sin dejar de ser por ello una mujer de carne y hueso.

Y sin embargo, Graves no escatima crueldad, incluyendo ---para mi pesar, yo las habría indultado--- la ejecución de las esclavas. Lo hace porque prefiere no darnos gato por liebre. La voz de Nausícaa es deliciosa, su personaje memorable, la gracia y el estilo de la historia son puro néctar, pero la matanza es igual de escalofriante que el original.

Y es que...

\emph{El clásico existe para que las generaciones posteriores lo reescriban, lo reimaginen, lo reinterpreten, pero nunca para que lo adulteren.}

El clásico existe para que cada generación descubra que la vida es corta, pero el arte es inmortal.
